% ============================================================================
\section{Foundations}
\label{sec:foundations}
% Source: assignments/01_foundations/index.html
% ============================================================================

% ----------------------------------------------------------------------------
\subsection{Optimization and probability}

\problempart{1a}{Weighted quadratic minimisation}
\question{Let $x_1,\dots,x_n\in\R$ and weights $w_1,\dots,w_n>0$. For
$f(\theta)=\sum_{i=1}^n w_i(\theta-x_i)^2$ ($\theta$ scalar), find the minimiser, show it is a minimum,
and explain what goes wrong if some $w_i$ are negative.}
\expect{Expression for the minimising $\theta$ and its derivation; a short argument that it is a minimum;
1--2 sentences on negative weights.}
\begin{solution}
$f$ is differentiable, so any minimiser is a critical point:
\[
f'(\theta)=\sum_{i=1}^n 2w_i(\theta-x_i)=2\theta\sum_i w_i-2\sum_i w_ix_i=0
\quad\Longrightarrow\quad
\theta^\star=\frac{\sum_{i=1}^n w_i x_i}{\sum_{i=1}^n w_i},
\]
the weighted average of the positions. It is a minimum because
$f''(\theta)=2\sum_i w_i>0$ (all $w_i>0$), so $f$ is strictly convex and its unique critical point is the
global minimum. (Division by $\sum_i w_i$ is safe because the sum is positive.)

If some $w_i$ are negative, two things can break: $\sum_i w_i$ may be $0$, so $\theta^\star$ is undefined,
and if $\sum_i w_i<0$ then $f''<0$, so the critical point is a \emph{maximum} and $f(\theta)\to-\infty$ as
$|\theta|\to\infty$, i.e.\ there is no minimum at all. Intuitively a negative weight ``rewards'' being far from
$x_i$, so the problem is no longer about staying close to the points.
\end{solution}

\problempart{1b}{Swapping min and sum}
\question{Let $f(\bvec x)=\min_{s\in[-1,1]}\sum_{i=1}^d s\,x_i$ and
$g(\bvec x)=\sum_{i=1}^d\min_{s_i\in[-1,1]} s_i x_i$ for $\bvec x\in\R^d$.
Which of $f\le g$, $f=g$, $f\ge g$ holds for all $\bvec x$? Prove it.}
\expect{A short (3--5 line) proof using mathematical notation.}
\begin{solution}
$f(\bvec x)\ge g(\bvec x)$ for all $\bvec x$.
Write both as minimisations of the same quantity $\sum_i s_i x_i$:
\[
f(\bvec x)=\min_{\substack{s_1,\dots,s_d\in[-1,1]\\ s_1=s_2=\dots=s_d}}\ \sum_{i=1}^d s_i x_i,
\qquad
g(\bvec x)=\min_{s_1,\dots,s_d\in[-1,1]}\ \sum_{i=1}^d s_i x_i .
\]
$f$ minimises over the subset of $[-1,1]^d$ where all coefficients are equal, while $g$ minimises the same
objective over the whole cube. Minimising over a larger set can only give a smaller (or equal) value, so
$g(\bvec x)\le f(\bvec x)$.
Concretely, $f(\bvec x)=-\bigl|\sum_i x_i\bigr|$ (take $s=-\operatorname{sign}\sum_i x_i$) and
$g(\bvec x)=-\sum_i|x_i|$ (take $s_i=-\operatorname{sign}x_i$), and the triangle inequality
$\bigl|\sum_i x_i\bigr|\le\sum_i|x_i|$ gives the same conclusion. The inequality is strict whenever the $x_i$ do
not all have the same sign, e.g.\ $\bvec x=(1,-1)$: $f=0>g=-2$.
\end{solution}

\problempart{1c}{Expected points of a dice game}
\question{Roll a fair six-sided die until a $1$ appears. Each $3$ wins $a$ points, each $6$ loses $b$
points, and $2,4,5$ change nothing. What is the expected number of points at the end, as a function of $a,b$?}
\expect{A recurrence for the value $V$, a short explanation of it, and the solved expression.}
\begin{solution}
Let $V$ be the expected number of points collected from now until the game stops. Each roll is independent,
so after any roll that does not stop the game we are in exactly the same situation and expect $V$ more points.
Conditioning on the next roll:
\[
V=\underbrace{\tfrac16\cdot 0}_{\text{roll }1\text{: stop}}
 +\underbrace{\tfrac16\,(a+V)}_{\text{roll }3}
 +\underbrace{\tfrac16\,(-b+V)}_{\text{roll }6}
 +\underbrace{\tfrac36\,V}_{\text{roll }2,4,5}
 =\frac{a-b}{6}+\frac56 V .
\]
Solving, $\tfrac16 V=\tfrac{a-b}{6}$, hence $\boxed{V=a-b}$.
(Sanity check: before the stopping $1$ appears, each of the other five faces is equally likely, so on
average there is exactly one $3$ and one $6$.)
\end{solution}

\problempart{1d}{Maximum likelihood for a coin}
\question{A coin with $P(\text{H})=p$, $0<p<1$, gives $\{\text{T},\text{H},\text{H},\text{H},\text{H}\}$, so
$L(p)=p^4(1-p)$. Which $p$ maximises $L$? Prove it is a maximum and give an intuitive interpretation.}
\expect{The maximiser with its derivation (e.g.\ via $\log L$), proof that it is a maximum, and a
1-sentence interpretation.}
\begin{solution}
$\log$ is strictly increasing, so $L$ and $\ell(p)=\log L(p)=4\log p+\log(1-p)$ have the same maximiser.
\[
\ell'(p)=\frac4p-\frac1{1-p}=0\ \Longrightarrow\ 4(1-p)=p\ \Longrightarrow\ p^\star=\frac45 .
\]
It is a maximum: $\ell''(p)=-\dfrac4{p^2}-\dfrac1{(1-p)^2}<0$ for every $p\in(0,1)$, so $\ell$ is strictly
concave and $p^\star$ is its unique global maximiser (equivalently, $L(p)>0$ on $(0,1)$ while
$L(p)\to0$ at both ends, so the unique interior critical point must be the maximum).

Interpretation: $p^\star=4/5$ is simply the empirical frequency of heads in the data (4 heads out of 5 flips);
maximum likelihood says the most plausible coin is the one whose bias matches what we observed.
\end{solution}

\problempart{1e}{Conditional probability}
\question{Events $A,B$ satisfy $P(A\mid B)=P(B\mid A)$, $P(A\cup B)=\tfrac13$ and $P(A\cap B)>0$.
Prove that $P(A)>\tfrac16$.}
\expect{A short ($\sim$5 line) proof.}
\begin{solution}
By the definition of conditional probability,
$P(A\mid B)=\dfrac{P(A\cap B)}{P(B)}$ and $P(B\mid A)=\dfrac{P(A\cap B)}{P(A)}$.
Setting them equal and using $P(A\cap B)>0$ (so we may divide by it) gives $P(A)=P(B)$.
By inclusion--exclusion,
\[
\tfrac13=P(A\cup B)=P(A)+P(B)-P(A\cap B)=2P(A)-P(A\cap B),
\]
so
\[
P(A)=\frac{\tfrac13+P(A\cap B)}{2}>\frac{\tfrac13+0}{2}=\frac16 ,
\]
where the inequality is strict because $P(A\cap B)>0$.
\end{solution}

\problempart{1f}{Gradient}
\question{For $\bvec w\in\R^d$, constants $\bvec a_i,\bvec b_j\in\R^d$, $\lambda\in\R$ and $n\in\mathbb{N}$,
\[
f(\bvec w)=\Bigg(\sum_{i=1}^n\sum_{j=1}^n(\bvec a_i^\top\bvec w-\bvec b_j^\top\bvec w)^2\Bigg)
+\frac{\lambda}{2}\norm{\bvec w}^2 .
\]
Compute $\nabla f(\bvec w)$ (no summations over dimensions).}
\expect{A compact expression for the gradient and its derivation ($\sim$5 lines).}
\begin{solution}
Write each term as $(\bvec c_{ij}^\top\bvec w)^2$ with $\bvec c_{ij}=\bvec a_i-\bvec b_j$, since
$\bvec a_i^\top\bvec w-\bvec b_j^\top\bvec w=(\bvec a_i-\bvec b_j)^\top\bvec w$.
Two standard gradients: $\nabla_{\bvec w}(\bvec c^\top\bvec w)=\bvec c$ and
$\nabla_{\bvec w}\tfrac12\norm{\bvec w}^2=\bvec w$. By the chain rule,
$\nabla_{\bvec w}(\bvec c^\top\bvec w)^2=2(\bvec c^\top\bvec w)\,\bvec c$, so
\[
\nabla f(\bvec w)=2\sum_{i=1}^n\sum_{j=1}^n\bigl((\bvec a_i-\bvec b_j)^\top\bvec w\bigr)(\bvec a_i-\bvec b_j)
+\lambda\bvec w .
\]
Equivalently, with the $d\times d$ matrix $M=\sum_{i,j}(\bvec a_i-\bvec b_j)(\bvec a_i-\bvec b_j)^\top$,
$f(\bvec w)=\bvec w^\top M\bvec w+\tfrac{\lambda}{2}\norm{\bvec w}^2$ and
$\nabla f(\bvec w)=(2M+\lambda I)\,\bvec w$.
For $d=1$ this reduces to $f'(w)=2\sum_{i,j}(a_i-b_j)^2w+\lambda w$, as expected.
\end{solution}

% ----------------------------------------------------------------------------
\subsection{Complexity}

\problempart{2a}{Placing three rectangles}
\question{On an $n\times n$ grid of points, how many ways are there to place $3$ axis-aligned rectangles
whose corners are grid points (degenerate and overlapping rectangles allowed)? Answer as $O(n^c)$ or $O(c^n)$.}
\expect{A big-O bound and $\sim$2 sentences of reasoning.}
\begin{solution}
$O(n^{12})$. An axis-aligned rectangle is fully determined by the $x$-coordinates of its left and right
sides and the $y$-coordinates of its bottom and top sides, each of which is one of $n$ grid values (repeats
allowed, giving degenerate rectangles), so there are $\Theta(n^2)\cdot\Theta(n^2)=\Theta(n^4)$ rectangles.
The three rectangles are placed independently and may coincide, so the number of placements is
$\Theta(n^4)^3=\Theta(n^{12})$; whether the three are ordered or not only changes the constant ($3!$), not
the exponent.
\end{solution}

\problempart{2b}{Minimum-cost monotone path}
\question{On an $n\times 2n$ grid, go from $(1,1)$ to $(n,2n)$ using only down/right steps. Each point has
cost $c(i,j)$ (possibly negative, $O(1)$ to evaluate), and a path's cost is the sum over the points it visits.
Give the most efficient algorithm for the minimum path cost and its runtime.}
\expect{Algorithm description ($\sim$5 sentences), the big-O runtime, and why.}
\begin{solution}
Use dynamic programming. Let $C(i,j)$ be the minimum cost of a path from $(1,1)$ to $(i,j)$. Because every
step goes down or right, $(i,j)$ can only be entered from $(i-1,j)$ or $(i,j-1)$, so
\[
C(1,1)=c(1,1),\qquad
C(i,j)=c(i,j)+\min\bigl(C(i-1,j),\,C(i,j-1)\bigr),
\]
where a term with $i-1<1$ or $j-1<1$ is treated as $+\infty$ (so the first row and column just accumulate).
Fill the table row by row (any order in which $(i-1,j)$ and $(i,j-1)$ are computed before $(i,j)$ works), and
return $C(n,2n)$. Negative costs are not a problem: the move graph is acyclic, so no path can revisit a point
and the recurrence remains valid.

Runtime $O(n^2)$: there are $n\cdot 2n=2n^2$ cells and each costs $O(1)$ (one call to $c$ and one $\min$).
This is optimal, since any correct algorithm must look at every $c(i,j)$ (each point lies on some path and a
single very negative value could change the answer). Keeping only the previous row reduces memory to $O(n)$.
\end{solution}

% ----------------------------------------------------------------------------
\subsection{Ethical issue spotting}

For each scenario, either identify one concern from the NeurIPS Ethics Guidelines (General Ethical Conduct /
Potential Negative Social Impacts), copy it out and \uline{underline} it, and justify it, or argue that
there is no concern.

\expect{A 2--5 sentence paragraph per scenario.}

\problempart{3a}{Loan-default prediction}
\question{A loan-default model relies mainly on location, accepts urban and rejects rural applicants, and
has a significantly higher false-positive rate for Black and male applicants.}
\begin{solution}
This model raises a \emph{Discrimination / Bias and fairness} concern: it may
\uline{encode, contain, or potentially exacerbate bias against people of a certain gender, race,
sexuality, or who have other protected characteristics}. A higher false-positive rate for Black and male
applicants means that creditworthy people in those groups are denied loans more often than equally creditworthy
people elsewhere, which is a direct, material harm tied to protected characteristics. Relying mainly on
location makes this worse: residence is strongly correlated with race in many countries (historical
``redlining''), so location acts as a proxy that lets the model discriminate even without a race feature, and
the systematic rejection of rural applicants adds a geographic form of exclusion that can entrench existing
inequality in access to credit.
\end{solution}

\problempart{3b}{Stylometry}
\question{Highly accurate ML stylometry predicts the authors of anonymous texts (including code) and could
de-anonymise a wide range of writing.}
\begin{solution}
The concern here is \emph{Surveillance}: the technique could be used to
\uline{facilitate surveillance, profiling, or the de-anonymization of individuals}. Although the stated use
is historical and literary scholarship, the same model applied to blog posts, forum comments or source code
could unmask whistle-blowers, dissidents, journalists' sources or contributors to privacy-preserving software,
and expose them to retaliation; anonymity is often the only protection such authors have. Because the
capability is general and cheap to run at scale, the negative use is at least as likely as the benign one,
so the researchers should discuss this dual use and consider mitigations (e.g.\ restricted release of models,
research on authorship obfuscation as a countermeasure).
\end{solution}

\problempart{3c}{Celebrity face dataset}
\question{Millions of celebrity faces were scraped from Google Images, without permission and including
copyrighted images, to build facial recognition; copyrighted images are distributed as URLs plus bounding boxes.}
\begin{solution}
This violates the data-related rules of General Ethical Conduct, in particular \emph{Consent}: datasets
containing people \uline{should have been collected with the consent of the people depicted}, and the
\emph{Copyright and Fair Use} requirement that researchers \uline{respect the copyright of the data}. The
subjects never agreed to have their faces used to train a biometric system, and being famous does not waive
that interest. Distributing URLs plus bounding boxes instead of the images is a technicality that makes the
copyrighted data just as usable while avoiding responsibility for it. There is also a downstream
\emph{Surveillance} concern, since facial recognition trained this way can be applied to anyone, not only to
celebrities.
\end{solution}

\problempart{3d}{Plant-species identification}
\question{A model identifies plant species from a photo. It was trained on iNaturalist photos whose users
consented to research use, and it is only used inside the app.}
\begin{solution}
We do not think this scenario raises an ethical concern under the guidelines. The data were contributed by users
who explicitly \uline{consented to the use of their photos for research}, so the consent and privacy
requirements are satisfied; the subjects are plants rather than people, so there is no protected-group bias to
encode; and the deployment is narrow (identification inside the same app, at the user's request), which
limits the risk of misuse or negative side effects. The only residual risk we see is indirect: precise
photo locations of rare or protected species could help poachers, but that is a property of the public
observation data and the app's location policy, not of the identification model itself.
\end{solution}

% ----------------------------------------------------------------------------
\subsection{Programming}
\codepart{1.4a--1.4f}
Functions: \texttt{find\_alphabetically\_first\_word}, \texttt{euclidean\_distance},
\texttt{mutate\_sentences}, \texttt{sparse\_vector\_dot\_product}, \texttt{increment\_sparse\_vector},
\texttt{find\_nonsingleton\_words}. All six pass the local grader (8/8 points); the hidden tests run without
errors or timeouts.
